% =====================================================================================
\subsection{Without instructions: a recogniser, an audit, and a handicap}
\label{sec:audit}

The rule of this subsection reads no instruction. One classifier over all tasks decides the route
from a frozen feature of the whole input, so no task can be blocked and nothing has to be known
about how a task is worded. Its guarantees are of two kinds: (a) an \emph{audit}, an estimate and a
confidence bound, from held-out inputs, of the share of a task's answers that have changed, valid
for any recogniser that is fitted without those inputs; and (b) a \emph{handicap}, a limit, known
in advance, on what a single later task can take from an older one. The audit is ordinary held-out
estimation, and the handicap is a conformal rank argument with a union bound. What is new is only
this: by Lemma~\ref{lem:route}, both become statements about answers. We first state what no method without instructions can avoid.

\paragraph{The price of having no instruction.} Let two tasks $j < k$ have inputs with
probabilities $p_j(x)$ and $p_k(x)$ (inputs are texts, so there are countably many), and let
$\omega_i(x, y) \in [0, M]$ be the expected score of the answer $y$ to the input $x$ under task
$i \in \{j, k\}$, the expectation being over the reference answers of that task for $x$. Write
$\omega_i^*(x) = \sup_y \omega_i(x, y)$, and for a system that answers $y_T(x)$ at time $T$ let
$\ell_i(T) = \mathbb{E}_{X \sim p_i}\bigl[\omega_i^*(X) - \omega_i(X, y_T(X))\bigr]$ be what it loses on task
$i$ against the best possible answers.

\begin{proposition}[The price of having no instruction]
\label{prop:noinstructionprice}
For every system whose answer at time $T \ge k$, or the distribution of its answer if it answers at
random, depends on the input alone,
\[
  \ell_j(T) + \ell_k(T) \;\ge\; \Omega_{jk} \;:=\; \sum_x \min\bigl(p_j(x), p_k(x)\bigr)\, \gamma_{jk}(x),
\]
where
\[
  \gamma_{jk}(x) \;=\; \inf_y \Bigl[ \bigl(\omega_j^*(x) - \omega_j(x, y)\bigr) + \bigl(\omega_k^*(x) - \omega_k(x, y)\bigr) \Bigr] .
\]
\end{proposition}

\begin{proof}
Fix $x$, let $y = y_T(x)$, $a = \omega_j^*(x) - \omega_j(x, y) \ge 0$ and $b = \omega_k^*(x) - \omega_k(x, y) \ge 0$.
Then $p_j(x)\, a + p_k(x)\, b \ge \min(p_j(x), p_k(x))\,(a + b) \ge \min(p_j(x), p_k(x))\, \gamma_{jk}(x)$.
Sum over $x$. For a system that answers at random, $a$ and $b$ are averages over its answers, each
answer satisfies $a + b \ge \gamma_{jk}(x)$, and the same lines hold.
\end{proof}

Three things follow. (i) If no input belongs to both tasks, $\Omega_{jk} = 0$ and nothing is
forced. This is the case when every input begins with the instruction of its task and no
instruction is the beginning of another: it is the exact sense in which the instruction gives the
task away. (ii) If the score is $M$ for the right answer and $0$ otherwise, and every input that
belongs to both tasks has a right answer under each of them and the two differ, then
$\gamma_{jk} = M$ on those inputs and $\Omega_{jk}/M$ is the overlap of the two input
distributions, one minus their total-variation distance. (iii) The forgetting of task $j$ between
its arrival and time $T$ is $\ell_j(T) - \ell_j(j)$, so
\[
  \bigl(\ell_j(T) - \ell_j(j)\bigr) \;+\; \ell_k(T) \;\ge\; \Omega_{jk} - \ell_j(j) :
\]
what is forgotten of the older task plus what is not learned of the newer one is at least the
conflict $\Omega_{jk}$, minus what the older task never had. A switch with one frozen set of
adapters per task can choose how this sum is split. A rule under which the oldest claim wins
(the zones of Sections~\ref{sec:switch} and~\ref{sec:bound}, below) puts it on the newer task,
except for a share of inputs that is chosen in advance. The recogniser and its variant \handicap{}
put a measured, or a limited, part of it on the older task. Every other method, with shared weights or not, is bound by the same inequality, but
does not say how the sum is split. The proposition is elementary, and that the recognition of the
task limits learning without task labels is known \citep{kim2022theoretical}; we state it because
it gives the recogniser and its variant below their place, and because $\Omega_{jk}$ can be counted from below on a
benchmark (inputs that occur word for word in two tasks; Section~\ref{sec:nodefpublished}).

\paragraph{A joint recogniser.} Let $u(x) \in \mathbb{R}^d$ be a frozen feature of the whole
input (in our experiments the mean of the last-layer states of the frozen encoder of the base
model, without any adapter, standardised with the mean and standard deviation of a fixed pool of
instances of other tasks, Section~\ref{sec:protocol}), and
$\bar u(x) = (u(x), 1)$. When task $k$ arrives, its
training inputs are split at random into fit inputs $D_k^{\mathrm{f}}$ and $m$ held-out inputs
$c_1, \dots, c_m$. The fit inputs add one vector to what is stored and one matrix to a sum kept for
all tasks,
\[
  G_k \;=\; \frac{1}{|D_k^{\mathrm{f}}|} \sum_{x \in D_k^{\mathrm{f}}} \bar u(x)\, \bar u(x)^{\!\top}, \qquad
  v_k \;=\; \frac{1}{|D_k^{\mathrm{f}}|} \sum_{x \in D_k^{\mathrm{f}}} \bar u(x),
\]
so that every task has the same weight, and the recogniser after $T$ tasks is
\begin{equation}
  V_T \;=\; \Bigl(\, \sum_{k \le T} G_k + \varepsilon\, I \Bigr)^{-1} \bigl[\, v_1, \dots, v_T \,\bigr]
  \label{eq:recogniser}
\end{equation}
with a fixed $\varepsilon > 0$. After $T$ tasks the route
of an input is the task with the largest score, the older task in a tie,
\begin{equation}
  \rho_T(x) \;=\; \arg\max_{k \le T}\ \bigl(\bar u(x)^{\!\top} V_T\bigr)_k ,
  \label{eq:jointroute}
\end{equation}
and the answer is given as before by Eq.~\eqref{eq:answer}, with the adapters that existed when
task $\rho_T(x)$ was learned. $V_T$ is the ridge regression of the task labels on the features of
all tasks seen so far. Because the sums are only added, it is the same for every order of
arrival, no input has to be kept, and every task has a column of its own: no task is blocked. The
adapter of a task may be trained on all of its inputs, the held-out ones included; the recogniser
never sees those. A recogniser of this kind is known from class-incremental learning on frozen
features \citep{mcdonnell2023ranpac,zhuang2022acil}; what we add is what is said about its routes.

By Lemma~\ref{lem:route} an answer can change only if its route changes, for this rule as for
any other.

\begin{proposition}[An audit of forgetting from held-out inputs]
\label{prop:audit}
Fix a task $j$ and two times $T' > T \ge j$. Let the fit inputs of all tasks and the choice of the
later tasks be independent of the held-out inputs $c_1, \dots, c_m$ of task $j$ and of a further
input $X$ of task $j$, let the route at every time be fitted without them (as in
Eqs.~\eqref{eq:recogniser}--\eqref{eq:jointroute}), and let $c_1, \dots, c_m, X$ be exchangeable
given everything else. Let $H$ be the number of held-out inputs with
$\rho_{T'}(c_i) \ne \rho_T(c_i)$. Then
\[
  \Pr\bigl[\, y_{T'}(X) \ne y_T(X) \,\bigr] \;\le\; \Pr\bigl[\, \rho_{T'}(X) \ne \rho_T(X) \,\bigr]
  \;=\; \mathbb{E}\Bigl[\frac{H}{m}\Bigr] .
\]
If the inputs are independent draws, then with probability at least $1 - \delta$ over the held-out
inputs, the share of the inputs of task $j$ whose route differs between $T$ and $T'$ is at most
the one-sided Clopper--Pearson upper bound for $H$ successes in $m$ trials at confidence
$1 - \delta$. Both statements also hold when ``differs between $T$ and $T'$'' is replaced by ``is
not the same at all times from $T$ to $T'$'', with $H$ counted in the same way.
\end{proposition}

\begin{proof}
The first inequality holds input by input, as shown above. Let
$I(x) = \mathbf{1}[\rho_{T'}(x) \ne \rho_T(x)]$. Given everything but $c_1, \dots, c_m, X$ the two
routes are fixed functions, so the $m + 1$ numbers $I(c_1), \dots, I(c_m), I(X)$ are exchangeable
and $\Pr[I(X) = 1] = \mathbb{E}[I(c_i)] = \mathbb{E}[H]/m$. For independent draws, $H$ is, given
the routes, binomial with $m$ trials and success probability equal to the share in question, and
the second statement is the definition of the Clopper--Pearson bound. The last sentence follows
with $I(x)$ the indicator that the route of $x$ is not constant from $T$ to $T'$.
\end{proof}

So $H/m$ is an unbiased estimate of the share of changed routes, and the confidence bound is the
guarantee. The condition on independence is not empty: the later tasks may not be chosen by
looking at the held-out inputs, nor at anything computed from them, for example the answers of the
adapters of task $j$, which may have been trained on them. The bound is a statement for one task
and one pair of times. Read at every arrival and for all tasks, the bounds do not hold all at
once at the same confidence; for that the level has to be divided by the number of readings (the
``at all times'' form covers every time in between for one task with a single reading). The audit
needs no pass through the language model: only the features of the held-out inputs are kept. It
can be read before the adapter of a new task is trained, because the recogniser needs only
features: a system knows which older tasks a newcomer disturbs, and by how much, before it learns
the newcomer. The audit says nothing in advance, and it is no longer valid for held-out inputs
that were used to make a decision about the routes. The variant below uses them in exactly this way,
and comes with a bound in advance instead.

\paragraph{A handicap for the newcomer.} When task $k \ge 2$ arrives, compute $V_k$ by
Eq.~\eqref{eq:recogniser} and, for any input, the margin of the newcomer over the best older task,
$\xi_k(x) = (\bar u(x)^{\!\top} V_k)_k - \max_{i < k} (\bar u(x)^{\!\top} V_k)_i$. Fix an integer
$\nu \ge 0$. For each older task $j < k$, take the $(\nu+1)$-th largest value of $\xi_k$ among the
held-out inputs of task $j$ ($-\infty$ if task $j$ has at most $\nu$ of them). Let $\eta_k$ be the
largest of these values, or zero if that is larger. The newcomer \emph{takes} an input if
$\xi_k(x) > \eta_k$. The route is the newest task that took the input:
\begin{equation}
  \rho_T(x) \;=\; \max\bigl\{\, k \le T : \ k = 1 \ \text{ or } \ \xi_k(x) > \eta_k \,\bigr\} .
  \label{eq:handicap}
\end{equation}
The matrices $V_k$ and the thresholds $\eta_k$ are stored when task $k$ arrives and never changed.
So a newcomer takes an input only if the joint recogniser prefers it to every older task, and it
takes at most $\nu$ of the held-out inputs of each older task.

\begin{theorem}[A bound in advance without zones]
\label{thm:handicap}
Let the route be set by Eq.~\eqref{eq:handicap}, under the conditions of
Lemma~\ref{lem:route}. Let $X$ be an input of task $j$ that is exchangeable
with the $m$ held-out inputs $c_1, \dots, c_m$ of task $j$ given everything else, and let the
later tasks and their fit inputs be independent of $(c_1, \dots, c_m, X)$, for example fixed in
advance. If the sequence ends with task $T_{\mathrm{end}}$, then
the probability that the answer to $X$ changes at any time after task $j$ has arrived is at most
the probability that its route does, and
\[
  \Pr\bigl[\ \text{there are } j \le T < T' \le T_{\mathrm{end}} \text{ with } \rho_{T'}(X) \ne \rho_T(X) \ \bigr]
  \;\le\; \frac{(T_{\mathrm{end}} - j)\,(\nu+1)}{m+1} .
\]
\end{theorem}

\begin{proof}
By Eq.~\eqref{eq:handicap} the route of an input changes at the arrival of task $k$ only if the
newcomer takes it, that is, only if $\xi_k(X) > \eta_k$, and it changes at no other time, because
the margins and thresholds of the earlier tasks are frozen. The function $\xi_k$ is fitted on fit
inputs only, so given everything else the $m + 1$ numbers
$\xi_k(c_1), \dots, \xi_k(c_m), \xi_k(X)$ are exchangeable. The threshold $\eta_k$ depends on the
held-out inputs, but all that is used is that it is at least the $(\nu+1)$-th largest of the
first $m$ of these numbers. Among $m + 1$ exchangeable numbers, the last one is strictly larger
than the $(\nu+1)$-th largest of the others with probability at most $(\nu+1)/(m+1)$, ties
included. A union bound over the $T_{\mathrm{end}} - j$ later tasks gives the claim.
\end{proof}

\begin{remark}[What these statements do not say]
\label{rem:auditlimits}
\emph{(i) They bound change, not quality:} an input that the
recogniser gives to the wrong task from the start does not change, and is answered by the wrong
adapters. How many inputs go to their own task is measured, not guaranteed.
\emph{(ii) The price of the handicap is paid by the newcomer} and is seen when it arrives: the
share of its own held-out inputs that it does not take. With $\nu = 0$ one held-out input of an
older task that looks like the newcomer can be enough to keep most of the newcomer out.
\emph{(iii) The bound in advance grows with the number of later tasks and falls with the number of
held-out inputs.} For $m = 200$, $\nu = 2$ and $14$ later tasks it is $21\%$; what we measured on
test inputs was far smaller (Section~\ref{sec:recogniser}), because most later tasks take nothing
from most older ones, which the bound cannot know. A tighter bound in advance needs more held-out
inputs. They are not lost for training: the adapter may learn from them. Under the handicap the
audit of Proposition~\ref{prop:audit} is not valid for the same held-out inputs, because they set
the thresholds; a second held-out set would be needed for both.
\emph{(iv) All statements need test inputs that are like the training inputs of the task}
(exchangeability), and later tasks that were not chosen by looking at the held-out inputs or at
anything computed from them. The probabilities are over the held-out inputs and the test input
together; for one given split the share can be larger (compare
Remark~\ref{rem:boundlimits}(ii)).
\emph{(v) Nothing here reads an instruction.} When inputs do carry one, the feature can be computed
from the example alone while the adapters read the whole input. For a fixed input the answer then
still depends on the route only, and the route does not depend on how the instruction is worded.
So the audit and the bound hold for every wording of the instruction, and also when there is no instruction; the
exchangeability is then a statement about the examples. A system that receives one string has to
find where the example begins; the benchmarks give the two parts separately. A first test with trained answers is in Section~\ref{sec:exampleroutes}.
\end{remark}
